% =====================================================================
% Module 04 — Collapse-Surface Stratification
% Academic-Galactic v1.2
% =====================================================================

\section{Collapse-Surface Stratification}
\label{sec:CollapseSurfaceStratification}

\subsection{Overview}
Collapse-surface stratification defines how the topology of the Soft
80-dimensional manifold $\mathcal{M}^{80}$ evolves under curvature envelopes
$\Phi_7$, holonomy transport $\Phi_8$, and topology-modulation drift
$\Phi_8^t$. Collapse surfaces represent stratified submanifolds where geometric
structure changes discontinuously or semi-continuously.

This module integrates:
\begin{itemize}
    \item curvature envelopes (Module 02),
    \item holonomy transport and drift (Module 02),
    \item spectral continuity (Module 03),
    \item filtration structure (Module 01).
\end{itemize}

Modern curvature-discretization research
\citeClassA{LiPengDong2026,GutierrezUle2026}
and spectral curvature foundations
\citeClassA{SpectralCurvature2026}
provide the mathematical grounding.

\subsection{Collapse Surfaces}
\begin{definition}[Collapse Surface]
A collapse surface at altitude $k$ is a stratified submanifold
\[
C_k \subseteq \mathcal{F}_k
\]
where curvature envelopes $\Phi_7$ induce topology modulation.
\end{definition}

Collapse surfaces satisfy:
\[
\Phi_7|_{C_k} \neq 0,
\qquad
\Phi_8^t(C_k) \subseteq C_k,
\]
ensuring stability under drift.

\subsection{Topology Modulation}
Topology modulation is induced by the drift operator $\Phi_8^t$:
\[
\Phi_8^t : \mathcal{F}_k \to \mathcal{F}_k.
\]

Topology modulation causes:
\begin{itemize}
    \item stratified surface evolution,
    \item collapse-surface deformation,
    \item curvature-driven topology change.
\end{itemize}

This behavior is consistent with holonomy-curvature interaction
\citeClassA{BlackholeHolonomy2026}.

\subsection{Stratification Hierarchy}
Collapse surfaces form a stratification hierarchy:
\[
C_0 \subseteq C_1 \subseteq \cdots \subseteq C_{80}.
\]

Each $C_k$ is compatible with the filtration embedding:
\[
i_k(C_k) \subseteq C_{k+1}.
\]

This ensures altitude-preserving topology modulation.

\subsection{Discrete-to-Continuous Behavior}
Discrete curvature models \citeClassA{LiPengDong2026} converge to continuous
curvature envelopes $\Phi_7$:
\[
\lim_{n \to \infty} \Phi_7^{(n)} = \Phi_7.
\]

This convergence ensures:
\begin{itemize}
    \item stable collapse-surface evolution,
    \item consistent topology modulation,
    \item compatibility with spectral curvature (Module 03).
\end{itemize}

\subsection{Spectral Interaction}
Spectral curvature interacts with collapse surfaces:
\[
\kappa_{\mathrm{spec}}(C_k)
=
\left\| \frac{d}{dt} \Phi_5(\gamma(t)) \right\|
\quad
\text{for curves } \gamma \subseteq C_k.
\]

Spectral invariants remain stable under collapse-surface drift:
\[
I_{\mathrm{spec}}(C_k) = I_{\mathrm{spec}}(\Phi_8^t(C_k)).
\]

\subsection{Holonomy-Induced Collapse}
Holonomy transport induces collapse when curvature envelopes exceed a threshold:
\[
\|\Phi_7(p)\| > \tau_k \quad \Rightarrow \quad p \in C_k.
\]

Thresholds $\tau_k$ are altitude-dependent and compatible with the operator
algebra (Module 05).

\subsection{Summary}
Collapse-surface stratification defines topology modulation, collapse-surface
evolution, spectral interaction, and holonomy-induced deformation. These
structures form the topological backbone of the A7 operator altitude and
interact directly with curvature envelopes, spectral continuity, and the
composite operator algebra.

% =====================================================================
% End Module 04
% =====================================================================
